\section{Key Management}


\subsection{Key Hierarchy}

\begin{figure}[H]
\centering
\begin{tikzpicture}[
    node distance=1.2cm,
    box/.style={rectangle, draw, minimum width=3cm, minimum height=0.7cm,
                align=center, font=\footnotesize},
    arrow/.style={->, thick}
]
\node[box] (master) {Master Key\\(Hanzo KMS)};
\node[box, below left=1.2cm and 0.5cm of master] (iam-key) {IAM Identity\\(age)};
\node[box, below=1.2cm of master] (ats-key) {ATS Identity\\(age)};
\node[box, below right=1.2cm and 0.5cm of master] (kms-key) {KMS Identity\\(age)};
\node[box, below=1.2cm of iam-key] (iam-recip) {IAM Recipients};
\node[box, below=1.2cm of ats-key] (ats-recip) {ATS Recipients};
\node[box, below=1.2cm of kms-key] (kms-recip) {KMS Recipients};

\draw[arrow] (master) -- (iam-key) node[midway,left,font=\scriptsize] {HKDF};
\draw[arrow] (master) -- (ats-key) node[midway,left,font=\scriptsize] {HKDF};
\draw[arrow] (master) -- (kms-key) node[midway,right,font=\scriptsize] {HKDF};
\draw[arrow] (iam-key) -- (iam-recip);
\draw[arrow] (ats-key) -- (ats-recip);
\draw[arrow] (kms-key) -- (kms-recip);
\end{tikzpicture}
\caption{Key derivation hierarchy. Per-service age identities are derived
from the master key via HKDF.}
\label{fig:keys}
\end{figure}

\begin{enumerate}[nosep]
    \item \textbf{Master encryption key.} Stored in Hanzo KMS (HIP-0027).
    This is the root of trust for all replicate encryption. It is a 256-bit
    symmetric key used as HKDF input keying material.

    \item \textbf{Per-service age keypairs.} Derived from the master key via
    HKDF with the service name as the info parameter:
    \texttt{HKDF-SHA256(master, salt="hanzo-replicate", info=service\_name)}.
    The derived 32 bytes are used as the X25519 private key (age identity).
    The corresponding public key is the age recipient.

    \item \textbf{Recipients file.} Contains the public keys (recipients) of
    all parties authorized to decrypt the replica. For normal operation, this
    is the service's own recipient. For disaster recovery, this includes a
    recovery recipient held offline.
\end{enumerate}

\subsection{Key Rotation}

Key rotation proceeds without downtime:

\begin{enumerate}[nosep]
    \item Generate a new age keypair (derived from a new HKDF info or a
    rotated master key).
    \item Add the new recipient to the recipients file. The old recipient
    remains.
    \item Restart the \texttt{replicate} sidecar. New WAL segments and
    snapshots are encrypted to both old and new recipients.
    \item After the retention period expires (72 hours for WAL, 30 days for
    snapshots), all surviving replicas are encrypted to the new recipient.
    \item Remove the old recipient from the recipients file.
    \item The old identity is retained in KMS for historical reads. It is
    marked as \texttt{rotated} but never deleted.
\end{enumerate}

\subsection{Emergency Recovery}

Full restore requires two keys: the master key from Hanzo KMS and the
per-service age identity derived from it. The recovery procedure:

\begin{enumerate}[nosep]
    \item Retrieve the master key from KMS (requires KMS admin
    authentication).
    \item Derive the service identity:
    \texttt{age-keygen -o identity.txt} from the HKDF output.
    \item Run \texttt{replicate restore -config replicate.yml
    /data/base/data.db}.
    \item The restored database is encrypted with sqlcipher. The sqlcipher key
    is also in KMS.
\end{enumerate}

If KMS itself is unavailable (catastrophic failure), the KMS database is
restored first using the offline recovery identity (stored in a physical
safe or HSM).

% ============================================================================
% 6. DEPLOYMENT
% ============================================================================
